\subsection{Задача Tex to SQL}

% TODO: добаить ссылки на статьи из NotebookLM

\texttt{Text-to-SQL} — это задача, которая заключается в автоматическом переводе вопросов или запросов на естественном языке (NL) в запросы на языке структурированных запросов (SQL), которые могут быть выполнены на заданной реляционной базе данных для получения ответа

В задаче Text to SQL имеется несколько исторически сложившихся аспектов:

\begin{enumerate}
    \item \textbf{Преобразование языка:} Text-to-SQL функционирует как задача преобразования последовательности в последовательность (sequence-to-sequence problem);
    \item \textbf{Входные данные:} На вход подаётся вопрос на естественном языке (Q) и схема базы данных (D или S), которая включает информацию о таблицах, столбцах, их типах данных, первичных и внешних ключах, а также другие метаданные;
    \item \textbf{Цель:} сгенерировать точный и допустимый SQL-запрос ($\hat{Y}$), который корректно выражает намерение пользователя и может быть выполнен для получения требуемого результата из базы данных;
    \item \textbf{Упрощение доступа:} Системы Text-to-SQL значительно упрощают доступ к данным для пользователей, не обладающих техническими навыками или специализированными знаниями SQL, позволяя им взаимодействовать с базами данных с помощью обычного языка
\end{enumerate}

Подходы к решению данной задачи постоянно менялись, начиная от ранних методов, таких как \textbf{rule-based methods}, которые позже заменили подходы основанные на глубоком обучении, включая архитектуры типа «кодировщик-декодировщик» (encoder-decoder), такие как \textbf{LSTM} и \textbf{трансформеры}.

В последнее время доминирующем направлением стало использование больших языковых моделей (LLM), таких как серии GPT, благодаря их продвинутым возможностям генерации текста и рассуждений. Поэтому в контексте LLM основная проблема Text-to-SQL сводится к разработке оптимального промпта (prompt engineering), чтобы модель сгенерировала правильный SQL-запрос.



% ============================================================
% SUBSUBSECTION: Формулировка задачи text-to-SQL
% ============================================================
\subsubsection{Формулировка задачи text-to-SQL}

Задача: Обработка вопроса на \textbf{NL} в запрос \textbf{SQL}.

Входные и выходные данные:

\begin{table}[htbp]
    \centering
    \footnotesize % Еще меньше шрифт
    \caption{Компоненты системы Text-to-SQL}
    \begin{tabular}{|p{2.6cm}|p{4.4cm}|p{2.5cm}|p{4cm}|}\hline
    \textbf{Компонент} & \textbf{Input} & \textbf{Output} & \textbf{Форма} \\ \hline
    
    Основной запрос & 
    Вопрос на ест. языке (Q) & 
    SQL-запрос ($\hat{Y}$) & 
    Прямая генерация \\ \hline
    
    Схема БД & 
    Структура БД (таблицы, столбцы, ключи) & 
    — & 
    — \\ \hline
    
    Промпт & 
    Инструкция задачи & 
    — & 
    Цепочка рассуждений (CoT) \\ \hline
    
    Контекст & 
    Примеры Q/SQL (few-shot) & 
    — & 
    Самокоррекция \\ \hline
    
    Дополнения & 
    Правила/ограничения & 
    — & 
    CTEs (структурир. SQL) \\ \hline
    
    \end{tabular}
    \label{tab:sql-components}
\end{table}



% ============================================================
% SUBSUBSECTION: Архитектурные подходы и современные методы Text-to-SQL
% ============================================================
\subsubsection{Архитектурные подходы и современные методы Text-to-SQL}

В современных системах \texttt{Text-to-SQL}, основанных на LLM, используется несколько ключевых архитектурных подходов, которые часто комбинируются для повышения точности и эффективности


\textbf{A. Промпт-инжиниринг (Prompt Engineering)}

Это основной подход, который лежит в основе LLM-based Text-to-SQL и не требует дополнительного обучения модели. Он включает в себя:

\begin{enumerate}

    \item Представление вопроса (Question Representation) Это то, как схема базы данных и вопрос представляются LLM. Рекомендуемые форматы:
    \begin{itemize}
        % TODO: исправить формулу
        \item Code Representation Prompt (CR): Представление структурированных знаний с использованием синтаксиса SQL (например, операторов CREATE TABLE), что включает полную информацию о базе данных, включая первичные и внешние ключи (Foreign Keys, FK). Этот подход часто оказывается предпочтительным, поскольку стимулирует способность LLM к работе с кодом
        \item OpenAI Demostration Prompt (OD): Использование специальных символов (например, \#) для комментирования схемы и добавление конкретных инструкций, например, «Complete sqlite SQL query only and with no explanation» (что доказало свою эффективность)
    \end{itemize}

    \item Обучение в контексте (In-Context Learning / Few-Shot) Предоставление модели нескольких примеров для изучения шаблонов преобразования. Ключевые аспекты этого метода:
    \begin{itemize}
        \item Выбор примеров (Example Selection): Важно выбирать наиболее полезные примеры (Q). Методология DAIL Selection (DAIL) показала высокую эффективность, поскольку она учитывает сходство как самого вопроса, так и предполагаемого SQL-запроса (скелета запроса), что важно для обучения LLM сопоставлению вопроса с запросом
        \item Организация примеров (Example Organization): Для баланса между качеством информации и количеством примеров (что ограничено лимитом токенов) была предложена DAIL Organization (DAIL). Этот метод представляет пары «вопрос-SQL» без повторного включения полной схемы базы данных в каждый пример, что повышает эффективность использования токенов
    \end{itemize}

    \item Усовершенствование рассуждений (Reasoning Enhancement) Эти методы направлены на повышение способности LLM обрабатывать сложные запросы
    \begin{itemize}
        \item Chain of Thought (CoT): Побуждение модели к генерации промежуточных шагов рассуждения (например, с помощью фразы "Let's think step by step") для разбиения сложной задачи на более простые, управляемые этапы
        \item Query Decomposition (QDecomp) / Least-to-Most: Разделение сложного NL-запроса на последовательность логических шагов или подзапросов
        \item ACT-SQL (Automatic CoT): Метод, автоматически создающий примеры CoT для обучения моделей, что делает процесс более масштабируемым
        \item Self-Consistency: Генерация нескольких потенциальных SQL-запросов и выбор наиболее согласованного из них (например, с использованием механизма голосования), что повышает надёжность
    \end{itemize}

\end{enumerate}


\textbf{B. Тонкая настройка (Supervised Fine-Tuning, SFT)}

Тонкая настройка LLM (или малых языковых моделей, SLM) на специализированных датасетах Text-to-SQL (например, Spider, BIRD) — это альтернативный или дополнительный подход, который значительно повышает производительность, особенно для открытых моделей

\begin{enumerate}
    \item \textbf{Для открытых LLM/SLM:} Тонкая настройка необходима для повышения точности и конкурентоспособности. Например, модель SQL\_DS\_100M\_ft, основанная на открытой архитектуре CodeS, после контролируемой тонкой настройки достигла высокой точности (до $72.2\%$ EX на тестах BIRD) при низких вычислительных затратах.
    
    \item \textbf{Типы настройки:} Используется как полная тонкая настройка (Full-Parameter Fine-Tuning) для достижения максимальной точности в специфических задачах, так и параметрически эффективная тонкая настройка (Parameter-Efficient Fine-Tuning), которая экономит ресурсы, настраивая только часть параметров.
    
    \item \textbf{Ограничение:} После тонкой настройки LLM могут терять способность к обучению в контексте (in-context learning).
\end{enumerate}

\textbf{C. Архитектуры ИИ-Агентов (LLM Agent Methods)}

Рассмотрим современные многоэтапные архитектуры, которые сочетают промпт-инжиниринг с механизмами самокоррекции и внешними инструментами:

\begin{enumerate}
    \item \texttt{DIN-SQL}: Эта архитектура декомпозирует задачу Text-to-SQL на четыре управляемых этапа: 
    \begin{itemize}
        \item связывание схемы (Schema Linking);
        \item классификация и декомпозиция запроса (например, на простой, не вложенный, вложенный)
        \item  генерация SQL и самокоррекция (Self-Correction).
    \end{itemize} Для сложных запросов (nested complex queries) используется разбиение проблемы на несколько шагов, при котором сначала решаются подзапросы, а затем генерируется окончательный ответ

    % TODO: Исправить формулы
    \item \texttt{DAIL-SQL}: Хотя это в первую очередь метод промпт-инжиниринга, он является интегрированным решением, которое использует эффективное представление кода (CR), интеллектуальный выбор примеров (DAIL) и эффективную организацию (DAIL). DAIL-SQL может быть дополнительно оснащён механизмом самосогласованности (Self-Consistency) для повышения точности.
    
    \item \texttt{MAC-SQL}: Это фреймворк для совместной работы нескольких агентов, который использует Decomposer (для пошагового разбиения сложных запросов), Selector (для ограничения базы данных релевантными данными) и Refiner (для исправления ошибочных SQL-запросов с использованием внешних инструментов выполнения SQL)
    
    \item \texttt{SQL-CRAFT}: Интегрирует интерактивные циклы исправления (Interactive Correction Loop, IC-Loop) на основе обратной связи от системы управления базами данных (СУБД) и Python-усиленное рассуждение (Python-enhanced reasoning) для логических операций и расчётов перед генерацией SQL
\end{enumerate}



% ============================================================
% SUBSUBSECTION: Оценка и Бенчмарки
% ============================================================
\subsubsection{Оценка и Бенчмарки}

Оценка и бенчмаркинг в области Text-to-SQL с использованием больших языковых моделей (LLM) фокусируются на нескольких ключевых метриках, которые оценивают точность, корректность выполнения и эффективность сгенерированных SQL-запросов.

\textbf{Основные показатели}

\begin{enumerate}
    \item Точность выполнения (Execution Accuracy, EX)
    \item Точность точного совпадения (Exact Matching Accuracy, EM)
    \item Эффективность использования токенов (Token Efficiency / Cost Efficiency)
    \item Показатель достоверной эффективности (Valid Efficiency Score, VES)
    \item Точность тестового набора (Test-suite Accuracy, TS)
\end{enumerate}


% TODO: подумать как сделать текст ниже поменьше 
Точность выполнения (EX) является основным и наиболее важным показателем для оценки эффективности LLM-based решений. EX сравнивает результат выполнения (output) сгенерированного SQL-запроса с результатом выполнения эталонного (ground truth) SQL-запроса на экземплярах базы данных.

Точность точного совпадения (EM) — это более строгая метрика, которая измеряет синтаксическую корректность сгенерированного запроса. EM требует, чтобы сгенерированный SQL-запрос в точности совпадал (включая ключевые слова) с эталонным SQL-запросом. Но из-за разнообразия синтаксиса, эталонный SQL-запрос может быть не единственным корректным вариантом.

Эффективность использования токенов (Token Efficiency / Cost Efficiency) Эта метрика не является \underline{мерой корректности}, но является критически важной для практического применения Text-to-SQL, особенно при использовании проприетарных LLM, таких как модели OpenAI, которые взимают плату за токены.

Показатель достоверной эффективности (Valid Efficiency Score, VES) — это комплексный показатель, который оценивает как корректность, так и эффективность выполнения сгенерированного запроса. Он учитывает достоверность запроса (успешность выполнения и соответствие результата эталону) и эффективность запроса (скорость выполнения по сравнению с эталонным запросом).

Точность тестового набора (Test-suite Accuracy, TS) - эта метрика разработана для оценки производительности моделей на разнообразном и сфокусированном наборе тестовых баз данных. TS измеряет, насколько хорошо модель предсказывает SQL-запросы, которые дают правильные результаты в различных сценариях баз данных, и стремится измерить строгий верхний предел семантической точности

\textbf{Бенчмарки}

Оценка LLM-based Text-to-SQL решений проводится на основе крупномасштабных кросс-доменных наборов данных, которые служат стандартами для сравнения:

\begin{enumerate}
    \item Spider: Это крупномасштабный кросс-доменный набор данных, который является ключевым для Text-to-SQL. Он содержит сложные SQL-запросы. Производительность на нем измеряется, главным образом, Точностью выполнения (EX), и на его лидерборде ведущие решения основаны на LLM.
    \item BIRD: Бенчмарк, ориентированный на крупномасштабные базы данных, который стремится устранить разрыв между академическими исследованиями и реальными приложениями. BIRD специально исследует проблемы, связанные с большими базами данных с «грязными» значениями, необходимостью вывода внешнего знания и оптимизацией эффективности выполнения SQL.
    \item Spider-Realistic: Более сложный вариант Spider, где вопросы были вручную переработаны для моделирования более реалистичных сценариев, при этом эталонные SQL-запросы остались неизменными.
\end{enumerate}





